\documentclass[title=Monsieur, name=Je-Suis-Pas-Fatigué, r=0.41, g=0.70, b=0.11]{../mrms}
\begin{document}

\coverpage{../img/m-je-suis-pas-fatigue.jpg}
\collectionpage
\titlepage
\txtpage{
    Monsieur Je-Suis-Pas-Fatigué était un bonhomme resplendissant
    d'énergie, et qui ne ressentait jamais de fatigue. Ainsi, il n'avait jamais besoin d'arrêter.

    Un nouveau projet? Facile! Partir à l'aventure? Immédiatement! Visiter ses amis? Toujours!

    Épatant. Époustouflant. Incroyable? Littéralement!
}
\imgpage{./img/page-01.jpg}
\txtpage{
    Lorsque monsieur Je-Suis-Pas-Fatigué passait une bonne soirée avec ses amis, et que
    la nuit arrivait, monsieur Je-Suis-Pas-Fatigué s'exclamait:

    -- Moi, je ne ressens aucune fatigue! Va pour toute la nuit!

    Étrangement, il ne se rappelait jamais de ce qu'il avait fait ensuite. Ou sinon il avait
    vécu de drôles d'aventures invraisemblables, où on vole et où on est très léger, ou très lourd.
}
\imgpage{./img/page-02.jpg}
\txtpage{
    Un après-midi, monsieur Je-Suis-Pas-Fatigué était chez son ami, monsieur Ça-Fait-La-Job. Enfin,
    originellement il était arrivé tôt le matin, mais le temps avait filé entre leurs doigts.

    Ensemble, ils étaient en train de construire un parc dans la cours de monsieur Ça-Fait-La-Job.
}
\imgpage{./img/page-03.jpg}
\txtpage{
    Cependant, il restait encore beaucoup de travail, et monsieur Ça-Fait-La-Job bâilla:

    -- On devrait arrêter pour un moment et faire une sieste. Je cogne des clous.

    -- Balivernes, répondit monsieur-Je-Suis-Pas-Fatigué. Je me sens encore plein d'énergie. Tu
    peux faire ta sieste, et moi je vais continuer.
}
\imgpage{./img/page-04.jpg}
\txtpage{
    Deux heures plus tard, quand monsieur Ça-Fait-La-Job finit sa sieste, il retourna voir ce que
    monsieur Je-Suis-Pas-Fatigué avait pu faire entretemps. Il lui sembla que les choses n'avaient
    pas vraiment bougé... Y compris monsieur Je-Suis-Pas-Fatigué!
}
\imgpage{./img/page-05.jpg}
\txtpage{
    Le lendemain, monsieur Je-Suis-Pas-Fatigué avait invité monsieur Arrêter-Quand-Je-Veux, et
    madame Fin-De-Phrase, afin qu'ils puissent faire une partie de Brun.

    C'était le jeu favoris de monsieur Arrêter-Quand-Je-Veux. Un drôle de jeu où tout était dans un
    ton de brun, et où rien ne faisait vraiment de sens.
}
\imgpage{./img/page-06.jpg}
\txtpage{
    Monsieur Je-Suis-Pas-Fatigué avait organisé toute une grande journée:

    Dès le matin, ils se rejoignaient pour créer leurs personnages et débuter la partie. Puis ils
    allaient jouer jusqu'au dîner, et après une pause pique-nique, ils allaient continuer jusqu'au
    souper, à quel moment ils allaient commander un festin! Finalement, monsieur Je-Suis-Pas-Fatigué
    leur avait proposé de continuer toute la nuit, afin qu'ils puissent finalement arrêter au moment
    d'aller voir le lever du soleil.
}
\imgpage{./img/page-07.jpg}
\txtpage{
    À l'heure du dîner, il était temps de faire une bonne pause. Ils se préparèrent
    d'excellents sandwiches avant de partir se promener un peu.

    Il y avait un grand parc proche de chez monsieur Je-Suis-Pas-Fatigué qui était parfait.
}
\imgpage{./img/page-08.jpg}
\txtpage{
    Arrivés au parc, ils s'installèrent à l'ombre d'un arbre, et commencèrent le repas.

    -- Ces sandwiches sont excellents, surtout le... commença madame Fin-De-Phrase, entre deux bouchées.

    -- Jambon? Compléta monsieur Arrêter-Quand-Je-Veux, sans lever les yeux.
}
\imgpage{./img/page-09.jpg}
\txtpage{
    -- Enfin bon, continua madame Fin-De-Phrase, merci encore pour... pour...

    Puis elle se tourna vers monsieur Je-Suis-Pas-Fatigué, qui s'était allongé. Mais il ne répondit
    pas! Hélas, finalement madame Fin-De-Phrase et monsieur Arrêter-Quand-Je-Veux passèrent plutôt
    l'après-midi au parc, vu que monsieur Je-Suis-Pas-Fatigué avait besoin de plus de temps.
}
\imgpage{./img/page-10.jpg}
\txtpage{
    À son réveil, monsieur Je-Suis-Pas-Fatigué était incrédule. Il regarda l'heure sur sa montre.

    -- On est en retard! On ne terminera jamais notre partie de Brun à ce rythme-là.

    -- Pas le choix, faudra qu'on retarde l'heure de fin prévue afin de pouvoir terminer la partie, répondit monsieur
    Arrêter-Quand-Je-Veux.

    -- Mais non, rétorqua madame Fin-De-Phrase, ce sera plus sage si... si...

    -- Monsieur Je-Suis-Pas-Fatigué compléta: On le fait en plusieurs fois?
}
\imgpage{./img/page-11.jpg}
\txtpage{
    Quelques temps plus tard, monsieur Je-Suis-Pas-Fatigué était chez madame Je-Veux-pas-Dormir pour
    l'aider avec le mécanisme de son pendule. Elle jurait que l'autre jour il avait reculé, lorsqu'elle
    s'était allongé pendant quelle faisait un casse-tête.

    Monsieur Je-Suis-Pas-Fatigué était un expert auto-proclamé des mécanismes.
}
\imgpage{./img/page-12.jpg}
\txtpage{
    Arrivé au soir, ils n'avaient toujours pas réussi à trouver le problème.

    -- Monsieur Je-Suis-Pas-Fatigué, vous avez les yeux fermés, annonça madame Je-Veux-Pas-Dormir.

    -- Ah bon?

    -- Oui, continua madame Je-Veux-Pas-Dormir. De toute évidence, il est temps de préparer du café!

    Et ils firent une pause pour aller dans la cuisine pour manger un morceau et boire un café.
}
\imgpage{./img/page-13.jpg}
\txtpage{
    Revigoré, monsieur Je-Suis-Pas-Fatigué pu retourner investiguer le pendule.

    Tout le mécanisme était pourtant en bon état: Le pendule pendait, les aiguilles tournaient, etc...

    Tic. Toc. Tic. Toc.

    Fermant les yeux pour mieux se concentrer, monsieur Je-Suis-Pas-Fatigué essayait de trouver la réponse.
}
\imgpage{./img/page-14.jpg}
\txtpage{
    -- Ohlala! Dit madame Je-Veux-Pas-Dormir, en le brusquant de sa concentration, votre réflexion
    fait de drôles de bruits!

    Monsieur Je-Suis-Pas-Fatigué ouvrit les yeux de peine et de misère.

    -- Donnez-moi une minute madame Je-Veux-Pas-Dormir, répondit-t-il. Je pense que je dois m'allonger.

    -- Mais non! Sinon on ne pourra pas réparer mon pendule, s'indigna madame Je-Veux-Pas-Dormir.

    Monsieur Je-Suis-Pas-Fatigué partit alors pour la salle de bain, un oreiller à la main...
}
\imgpage{./img/page-15.jpg}
\txtpage{
    On entendit un clic de serrure, un foump, et puis... un ronflement!
}
\legalpage{28 juin 2026}
\backpage

\end{document}
